\section{Research Goal, Questions, and Methods}
\label{chap:rqm}
In this project, the Goal, Question, Metrics (GQM) approach \citep{basili_software_1992} is used to define the software metrics. This methodology introduces a conceptual, operational, and quantitative level. At the conceptual level, a goal for what the project should achieve is defined. The operational level includes a set of research questions associated with the goal, indicating whether the goal is accomplished. The quantitative level defines a set of metrics that will help answer the research questions.

\subsection{Research Goal}
\label{sec:research_goal}
The research goal of this project is:
\textit{The goal of this project is to invent an alternative game concept with a socially beneficial effect that can be developed, explored, and tested.}

\subsection{Research Questions}
\label{sec:research_questions}
To evaluate the potential of our game concept, we have defined seven research questions (RQs) decomposing the research goal:

\textbf{RQ1}: \textit{How do game metrics affect the players’ engagement in our game concept?}\\
\textbf{RQ2}: \textit{How do game metrics affect the players’ enjoyment in our game concept?}\\
\textbf{RQ3}: \textit{How do game metrics affect the players’ motivation in our game concept?}\\
\textbf{RQ4}: \textit{How does social interaction affect the players’ engagement in our game concept?}\\
\textbf{RQ5}: \textit{How does social interaction affect the players’ enjoyment in our game concept?}\\
\textbf{RQ6}: \textit{How does social interaction affect the players’ motivation in our game concept?}\\
\textbf{RQ7}: \textit{How does the use of our game concept affect the players’ habits?}\\

\textbf{RQ1, RQ2, and RQ3 description - Game Metrics}\\
These research questions will investigate how all game metrics implemented in our application will affect the players’ engagement, enjoyment, and motivation. All game metrics used in the application will be explored, e.g., points, leaderboard, challenges, and achievements. Concrete usage data from the application and direct feedback from players will be used to analyze and answer the questions.

\textbf{RQ4, RQ5, and RQ6 description - Social Interaction}\\
This research question will investigate how social interaction between players will affect the players’ engagement, enjoyment, and motivation in our application. Different types of social interaction will be evaluated and compared against each other. Concrete usage data from the application and direct feedback from players will be used to analyze and answer the questions.

\textbf{RQ7 description - Habits}\\
This research question will examine how the players’ habits evolve as a result of testing our game. The study will compare the players’ habits within the game’s socially beneficial focus areas before and after the test period. Answers to this question will primarily be based on responses from the questionnaire answered by all test participants. Due to the limited time of this master project, measuring actual user habits data is not possible.

\subsection{Research Methods}
\label{sec:research_methods}
To answer the research questions, we have developed a game using game theory principles that can be tested on a group of users. The methods we have used to accomplish our research goal are a literature review, to get the insight needed to design and develop an alternative game, and triangulation of data to answer our research question supported by both qualitative and quantitative data.
Literature review
To get the insight needed to design and develop an alternative game, we reviewed relevant literature during our specialization project. The research areas we reviewed included game development, game flow, existing alternative games, game platforms, and available game technology. Insight from this literature study on existing games, relevant technologies, and game development concepts are summarized in the prestudy (see Part \ref{part:prestudy}). 

\citet{aveyard_doing_2014} states that a literature review is a “comprehensive study and interpretation of literature that relates to a particular topic.” The purpose of our review is to combine theories from several articles to give the reader and writer a clear academic understanding of this thesis and the research field of alternative games. This insight is also helpful to conduct further research in this area of study.

\textbf{Triangulation of Data} \\
Triangulation of data refers to using multiple methods or data sources in qualitative research to corroborate findings and enhance the validity \citep{oates_researching_2006}. We use \textit{method triangulation}, utilizing two or more data generation methods. This approach increases the reliability of the result because the data can be cross-validated. 

Our analysis is based on both quantitative and qualitative data generation methods:

\textbf{Quantitative Data} \\
We have chosen three methods for quantitative data:
\begin{itemize}
    \item Questionnaire answered by the participants after the testing period. The questions are based on our research questions’ topics, i.e., the player's enjoyment, motivation, engagement, and change in habits.
    \item Screen time usage from the participants' smartphones.
    \item In-game data generated by participants.
\end{itemize}

\textbf{Qualitative Data} \\
To verify the effects of the game, we also used qualitative methods to corroborate our findings:

\begin{itemize}
    \item Semi-structured interviews.
    \item Participant observations on in-game interactions.
\end{itemize}


Part \ref{part:research_methodology} elaborate on our research methodologies, such as crucial choices, execution, and issues.
